\subsection{Further Intuitions on The Denoising Effect of Overparameterization}
To provide further insights into the pruning benefits of overparameterization, consider a simple linear model (as in Def~\ref{def LGP}) with $n\geq p\geq s$, noise level $\sigma=0$ and identity covariance. Suppose our goal is estimating the coefficients $\bt^\st_{\Delta}$ for some fixed index set $\Delta\subset[p]$ with $|\Delta|=s$. For pruning, we can pick $\Delta$ to be the most salient/largest entries. If we solve the smaller regression problem over $\Delta$, $\bth(\Delta)$ will only provide a noisy estimate of $\bt^\st_{\Delta}$. The reason is that, the signal energy of the missing features $[p]-\Delta$ acts as a noise uncorrelated with the features in $\Delta$. Conversely, if we solve ERM with all features (the larger problem), we perfectly recover  $\bt^\st$ due to zero noise and invertibility ($n\geq p$). Then one can also perfectly estimate $\bt^\st_{\Delta}$. This simple argument, which is partly inspired by the missing feature setup in \cite{hastie2019surprises}, shows that solving the larger problem with more parameters can have a ``denoising-like effect" and perform better than the small problem. Our contribution obviously goes well beyond this discussion and theoretically characterizes the exact asymptotics, handles the general covariance model and all $(n,p,s)$ regimes, and also highlights the importance of the overparameterized regime $n\ll p$.
%\textcolor{red}{It also explains why retraining again with few features can potentially hurt the performance. This is because when features are uncorrelated (e.g.~diagonal covariance), missing features (due to the smaller problem size) act like additive uncorrelated noise.} 
\section{Understanding the Benefits of Retraining}\label{sec intu}


%Without losing generality, let $[s]$ be the largest entries of latent parameter $\bt$ in absolute value. 
\cmt{\noindent\textbf{Denoising effect of overparameterization:} Consider a simple linear model with noise level $\sigma=0$, $n\geq p\gg s$ and identity covariance. Pick index set $\Delta\subset[p]$ with $|\Delta|=s$. Suppose we wish to estimate the coefficients $\bt^\st_{\Delta}$. For pruning, we can pick $\Delta$ to be the most salient/largest entries. If we solve the smaller problem over $\Delta$, $\bth[\Delta]$ will only provide a noisy estimate of $\bt^\st_{\Delta}$. The reason is that, the signal energy of the missing features $[p]-\Delta$ act as noise uncorrelated with features over $\Delta$. Conversely, if we solve ERM with all features (the larger problem), we perfectly recover  $\bt^\st$ due to lack of noise and $n>p$ from which we perfectly estimate $\bt^\st_{\Delta}$. This simple argument, which is in similar spirit to \cite{}, shows that solving the larger problem with more parameters can have a denoising like effect and perform better than the small problem. It also explains why retraining again with few features can hurt the performance. Our contribution obviously goes well beyond this discussion and theoretically characterizes exact asymptotics and also the important overparameterized regime $n\ll p$.}


%\noindent\textbf{Why \& When Retraining Helps:} 
%So far we have a couple of interesting conclusions. First, solving overparameterized problem and then pruning can be better than pruning with optimal sparse features. Secondly, retraining may hurt the performance in theory however it helps the performance for neural nets and random features. This section aims to provide further insights into these.

%\CT{I think this should be Fig. 7?}{
On the one hand, the study of LGPs in \fx{Fig.~\ref{fig1} and Fig.~7 of SM} suggest that retraining can actually hurt the performance. On the other hand, in practice and in the RFR experiments of Fig.~\ref{figRF2}, retraining is crucial; compare the green $\bth^M$ and red $\bth^{RT}$ curves and see SM Section A for further DNN experiments. Here, we argue that the benefit of retraining is connected to the correlations between input features. Indeed, the covariance/Hessian matrices associated with RF and DNN regression are not diagonal (as was the case in Fig.~\ref{fig1}). To build intuition, imagine that only a single feature suffices to explain the label. If there are multiple other features that can similarly explain the label, the model prediction will be shared across these features. Then, pruning will lead to a biased estimate, which can be mitigated by retraining. The following lemma formalizes this intuition under an instructive setup, where the features are perfectly correlated.% and the pruning performance can be exactly characterized.


\cmt{
train$\rightarrow$prune With growing correlations, the This is in contrast to the diagonal covariance with orthogonal features where retraining may hurt. wheredemonstrate a simple scenario which shows that retraining can actually
%  i.e.~when the same observation can be explained by multiple features
 retraining is critical for training sparse neural networks \cite{frankle2019lottery}. In this section, we identify feature correlation as a critical factor that necessitates retraining. This is in contrast to our results in Section \ref{doubdec} which considers independent features.

%Perhaps surprisingly, our results so far show that retraining might hurt the performance since estimating the feature coefficients by fitting the important features alone can perform worse than fitting all features.  

In order to convey the point, let us focus on a simplistic scenario where covariance matrix is rank one and features are perfectly correlated with each other. In this case, the diagonal entries of the covariance matrix essentially dictates the form of the solution. %Suppose the covariance is simply 
}
%Let $(\la_i)_{i=1}^p$ be the diagonal entries of $\sqrt{\bSi}$. Draw $\y,\X$ according to Definition \ref{d model}. 
\begin{lemma} \label{lem rank one}Suppose $\Sc$ is drawn from an LGP$(\sigma,\bSi,\bts)$ as in Def.~\ref{def LGP} where $\text{rank}(\bSi)=1$ with $\bSi=\bla\bla^T$ for $\bla\in\R^p$. Define $\zeta=\Tb_s(\bla)^2/\tn{\bla}^2$. For magnitude and Hessian pruning ($P\in\{M,H\}$) and the associated retraining, we have the following excess risks with respect to $\bt^\st$%population minima
\begin{align}%\underbrace{\frac{\zeta^2\sigma^2}{p-2}}_{\text{Variance Term}}
&\E_{\Sc}[\Lc(\bth_s^P)]-\Lc(\bt^\st)={\frac{\zeta^2\sigma^2}{n-2}}+\underbrace{(1-\zeta)^2(\bla^T\bts)^2}_{\text{Error due to bias}}\\
&\E_{\Sc}[\Lc(\bth_s^{RT})]-\Lc(\bt^\st)={\sigma^2}/({n-2}).
\end{align}
\end{lemma}
The lemma reveals that pruning the model leads to a biased estimator of the label. Specifically, the bias coefficient $1-\zeta$ arises from the missing predictions of the pruned features (which correspond to the small coefficients of $|\bla|$). In contrast, regardless of $s$, retraining always results in an unbiased estimator with the exact same risk as the dense model which quickly decays in sample size $n$. The reason is that retraining enables the remaining features to account for the missing predictions. Here, this is accomplished perfectly, due to the fully correlated nature of the problem. {In particular, this is in contrast to the diagonal covariance (Fig.~\ref{fig1}), where the missing features act like uncorrelated noise during retraining.} \cmt{Our distributional predictions in the next section enable us to capture the benefits of retraining for arbitrary covariances going beyond this instructive example.}% / feature correlations
%While this lemma provides insights for a simplistic model, w
%. They do this perfectly due to as
% and they can do it well since features are perfectly correlated
\cmt{
\begin{proof}
Set $c^\st=\bla^T\bt^\st$. By definition, each input example $\x_i$ has the form $\x_i=g_i\bla$ and $y_i=g_ic^\st+\sigma z_i$. Set $\g=[g_1~\dots~g_n]^T$ and $\bar{\g}=\g/\tn{\g}$. Thus, we have $\X=\g\bla^T$ and $\y=\g\bla^T\bt^\st+\sigma\z$. Decompose $\z=\bar{\z}+\bar{\g}^T\z\bar{\g}$. The least-squares solution has the form $\bth=\hat{c}\bla/\tn{\bla}^2$ where
\begin{align}
&\hat{c}=\arg\min_{c}\tn{(c^\st-c)\g+\sigma\z}\implies\\
&\hat{c}=c^\st+\sigma\gamma\label{beta}.
\end{align}
where $\gamma=\frac{\bar{\g}^T\z}{\tn{\g}}$. Set $h$, $\epsilon\distas \Nn(0,1)$. Note that 
\begin{align}
\Lc(\bth)=\mathbb{E}[((c^\st-\hat{c})h+\sigma\epsilon)^2]
=(\gamma^2+1)\sigma^2
\end{align}
\begin{align}
\Lc(\bth)&=\mathbb{E}[(c^\st- \hat{c}\frac{\bla^T\ub}{\tn{\bla}^2})h+\sigma\epsilon)^2]\\
&=\mathbb{E}[(c^\st- \hat{c}\zeta)h+\sigma\epsilon)^2]\\
&=((1-\zeta)c^\st-\zeta\sigma\gamma)^2
\end{align}
Now let $|\lambda_{k_1}|\geq|\lambda_{k_2}|\geq...\geq|\lambda_{k_p}|$. Assume we solved ERM to get  $\hat{c}$ and $\bth=\hat{c}\bla$. Prune the trained model $\bth$ to $s$-sparse by keeping the largest entries. Set $\ub=\Tb_s(\bla)$ and $\zeta=\tn{\ub}^2/\tn{\bla}^2$. We get $\bth_s=\hat{c}\ub$.
\begin{align}
    \Lc(\tbh)&=\mathbb{E}[((\bla^T\bt-\ub^T\tbh)h+\sigma\epsilon)^2]\\
    &=\mathbb{E}[((\bla^T\bt-\sum_{i=1}^s\lambda_{k_i}^2\hat{c})h+\sigma\epsilon)^2]\\
    &=[(1-\zeta)\bla^T\bt-\zeta\sigma\gamma]^2+\sigma^2
\end{align}
\textbf{Case 1: }
Now, assume that we have already known the optimal entries and do pruning before training. Then the pruned model can be seen as pruning the inputs and then putting it through a non-sparse $s$-feature model. Let $\lambda_{t_1}\bt_{t_1}\geq\lambda_{t_2}\bt_{t_2}\geq...\geq\lambda_{t_s}\bt_{t_s}$, $\w=[\lambda_{t_1}~\lambda_{t_2}~...~\lambda_{t_s}]$ and $\tb=[\bt_{t_1}~...~\bt_{t_s}]$. Set data sample $\x_i'\distas \Nn(0,\bSi')\in\R^s$ where are $\bSi'=\w\w^T$. Following what done in \ref{beta}, similarly
\[
\tbh'=\theta\w\quad\text{where}\quad\theta=\frac{1}{\sum_{i=1}^s\lambda_{t_i}^2}(\w^T\tb+\frac{\sigma}{\tn{\g}^2}\g^T\z)
\]
\begin{align}
    \Lc(\tbh')&=\mathbb{E}[((\bla^T\bt-\sum_{i=1}^s\lambda_{t_i}^2\theta)h+\sigma\epsilon)^2]\\
    &=[(\bla^T\bt-\w^T\tb)-\sigma\gamma]^2+\sigma^2
\end{align}
\textbf{Case 2: }
Now retrain with the pruned entries $\lambda_{k_!}$, ..., $\lambda_{k_s}$. Then the pruned model can be seen as pruning the inputs and then putting it through a non-sparse $s$-feature model. Let $\tb=[\bt_{k_1}~...~\bt_{k_s}]$. Set data sample $\x_i'\distas \Nn(0,\bSi')\in\R^s$ where are $\bSi'=\ub\ub^T$. Following what done in \ref{beta}, similarly
\[
\tbh'=\theta\vct u\quad\text{where}\quad\theta=\frac{1}{\sum_{i=1}^s\lambda_{k_i}^2}(\vct u^T\tb+\frac{\sigma}{\tn{\g}^2}\g^T\z)
\]
\begin{align}
    \Lc(\tbh')&=\mathbb{E}[((\bla^T\bt-\sum_{i=1}^s\lambda_{k_i}^2\theta)h+\sigma\epsilon)^2]\\
    &=[(\bla^T\bt-\vct u^T\tb)-\sigma\gamma]^2+\sigma^2
\end{align}
\textbf{It is obvious to know that loss in case 1 is always smaller than it in case 2. Then for both case, }

Set $\bt^T\bSi\bt=\alpha^2$ and $\tb^T\bSi'\tb=\rho^2$. We can write $\Lc(\tbh)$ and $\Lc(\tbh')$ as
\[
\Lc(\tbh)=((1-\zeta)\alpha-\zeta\sigma\gamma)^2+\sigma^2\iff\Lc(\tbh')=(\alpha-\rho-\sigma\gamma)^2+\sigma^2
\]
Now consider a special case which satisfies $\Lc(\tbh)<\Lc(\tbh')$.
\[
0\leq(1-\zeta)\alpha-\zeta\sigma\gamma<\alpha-\rho-\sigma\gamma
\]
\[
\frac{\rho+\sigma\gamma}{\alpha+\sigma\gamma}<\zeta\leq\frac{\alpha}{\alpha+\sigma\gamma}
\]
\end{proof}}

\cmt{
\begin{proof} Let $\bSi=\bla\bla^T$. Each $\x_i$ has the form $\x_i=g_i\bla$ and $y_i=x_i\bt^T\bla+\sigma z_i$. Thus, we have $\X=\g\bla$ and $\y=\g\bla^T\bt+\sigma\z$. Let $\z=\bar{\z}+\frac{\g^T\z}{\tn{\g}^2}\g$. The least-squares solution has the form $\bth=\hat{c}\bla$ where
\[
\hat{c}=%\arg\min ()
\]
\end{proof}}